\clearpage
\VerdierPhysicalPageAnchor{0263}%
\begingroup
\small
\setlength{\parindent}{0pt}%
\setlength{\parskip}{0pt}%
\setlength{\emergencystretch}{2em}%
\newcommand{\VerdierTermRow}[3]{%
  \VerdierSemanticAnchor{#1}%
  \par\noindent\hangindent=1.5em\hangafter=1
  #2。\enspace\textup{#3。}\par\smallskip}

\VerdierTermRow{idxt-p263-r01}{复形的第一典范滤过}{第三章，4.1.2}
\VerdierTermRow{idxt-p263-r02}{函子关于一个复形的第一谱序列}{第三章，4.4.1、4.4.4}
\VerdierTermRow{idxt-p263-r03}{$G$-范畴与 $G'$-范畴的积}{第一章，1.3.11}
\VerdierTermRow{idxt-p263-r04}{拟同构}{第三章，1.1.3}
\VerdierTermRow{idxt-p263-r05}{符号规则}{第一章，1.6.6}
\VerdierTermRow{idxt-p263-r06}{$M$ 型右分解（相应地，左分解）}{第三章，1.1.4}
\VerdierTermRow{idxt-p263-r07}{Cartan--Eilenberg 内射分解}{第三章，4.6.5}
\VerdierTermRow{idxt-p263-r08}{有限长度的 Cartan--Eilenberg 内射分解}{第三章，4.6.5}
\VerdierTermRow{idxt-p263-r09}{$\Sigma$-扩张（函子 $\int_\varphi$）}{第一章，2.2.3}
\VerdierTermRow{idxt-p263-r10}{函子到复形同伦范畴的 $\Sigma$-扩张}{第一章，2.5.10}
\VerdierTermRow{idxt-p263-r11}{多加法函子的 $\Sigma$-扩张}{第一章，2.4.2}
\VerdierTermRow{idxt-p263-r12}{函子态射到复形同伦范畴的 $\Sigma$-扩张}{第一章，2.5.10}
\VerdierTermRow{idxt-p263-r13}{乘法系的饱和化}{第二章，2.2.5}
\VerdierTermRow{idxt-p263-r14}{态射的分裂}{第二章，1.2.8}
\VerdierTermRow{idxt-p263-r15}{两个同伦之和}{第一章，2.5.4}
\VerdierTermRow{idxt-p263-r16}{右（或左）$\mathcal B$-局部化子范畴}{第二章，2.3.5}
\VerdierTermRow{idxt-p263-r17}{阿贝尔范畴的导出范畴中有界（相应地，下有界、上有界）对象的子范畴}{第三章，1.2.4}
\VerdierTermRow{idxt-p263-r18}{三角子范畴}{第二章，2.1.5}
\VerdierTermRow{idxt-p263-r19}{饱和三角子范畴}{第二章，2.1.6}
\VerdierTermRow{idxt-p263-r20}{相联系的函子列}{第三章，3.3.3}
\VerdierTermRow{idxt-p263-r21}{与谱对象关联的谱序列}{第二章，4.3.3}
\VerdierTermRow{idxt-p263-r22}{收敛谱序列}{第二章，4.4.3}
\VerdierTermRow{idxt-p263-r23}{Cartan--Eilenberg 谱序列}{第三章，4.6.9}
\VerdierTermRow{idxt-p263-r24}{定常谱序列}{第二章，4.4.3}
\VerdierTermRow{idxt-p263-r25}{与三角子范畴关联的乘法系}{第二章，2.1.15}
\endgroup
